% Generated packaging derivative; source remains byte-preserved.



\noindent\textbf{Status and main limitation.} This is a paper proof checked in
separate agent reviews, not a Lean-checked or externally specialist-reviewed
result. The review records accompany this note; model diversity is not
guaranteed. It gives matching characterizations for unbiased and uniformly
small-bias \emph{single-coin} estimators, including all moment--work tradeoffs
in those classes. It does not prove a lower
bound for arbitrary population algorithms and does not refute the critical
minimax conjecture. The latter remains open in this attempt.

\section{Exact scope and strongest statement}
Fix $a=1/4$ and the probability law
\[
 \mu(dd)=\frac{2d}{a^2}\ind_{(0,a]}(d)\,dd.
\]
The choice of one fixed interior $a$ suffices for the obstruction. Let
$0<\rho\le1/4$ and
$\mathcal Q_\rho=\{q:0\le q\le1,\ \int q\,d\mu\le\rho\}$,
with all functions measurable. A fresh coin has hidden magnitude $D\sim\mu$,
hidden sign $S\mid D=d\sim\operatorname{Ber}(q(d))$, and iid Bernoulli flips
of parameter $1/2+(2S-1)D$ conditional on $(D,S)$.

A single-coin procedure $A$ may know $\mu,a,\rho$ and may use arbitrary
independent randomization. It observes the flips of this one coin only,
stops at a possibly unbounded observable stopping time
$T\in\mathbb N_0\cup\{\infty\}$, and returns a real
observable output $Y$. It may stop before the first flip. It cannot sample
additional genuine reservoir coins; free parameter-independent simulated
variables can be included in its seed. Neither $D$, $S$, nor $q$ is observed.
Assume that for every $q\in\mathcal Q_\rho$,
\[
 \E_qY=\theta(q):=\int q\,d\mu,
 \qquad
 V=\sup_{q\in\mathcal Q_\rho}\Var_qY<\infty,
 \qquad
 W=\sup_{q\in\mathcal Q_\rho}\E_qT<\infty.
\]
Write $M=\sup_q\E_qY^2$. This is finite because $M\le V+\rho^2$.
All suprema in this note are over the stated nonempty class.
Nontermination is recorded as $T=\infty$, so finite expected work implies
almost-sure termination under every admissible population. The output can be
defined to be zero on the resulting null nontermination event.

\begin{theorem}[Sharp critical single-coin product]\label{thm:main}
There exist finite universal constants $0<c<C$ such that for every
$0<\rho\le1/4$, with $L=1+\log(1/\rho)$,
\[
 c\rho L(1+\log L)
 \le \inf_A V(A)W(A)
 \le C\rho L(1+\log L).                                  \tag{1}
\]
The infimum ranges over exactly the single-coin procedures above and permits
dependence on $\rho$. The same assertion holds with $M(A)$ in place of $V(A)$.
The lower bound holds for every admissible procedure; no threshold family,
finite depth, or numerical grid restriction is imposed.
\end{theorem}

The next theorem strengthens the lower bound beyond exact unbiasedness while
preserving Theorem~\ref{thm:main} as a separate statement.

\begin{theorem}[Small relative bias does not remove the entropy loss]\label{thm:bias}
Keep the same fixed law and single-coin observation rules. Replace exact
unbiasedness by
\[
 \sup_{q\in\mathcal Q_\rho}|\E_qY-\theta(q)|\le\rho/64,
\]
and retain separate finite uniform variance $V$ and expected work $W$.
There are universal $0<c<C<\infty$ such that, for every $0<\rho\le1/4$,
\[
 c\rho L(1+\log L)
 \le\inf_A V(A)W(A)
 \le C\rho L(1+\log L).                                 \tag{1b}
\]
The same assertion holds with $M(A)=\sup_q\E_qY^2$ in place of $V(A)$.
The infimum in (1b) includes every single-coin policy meeting this bias bound,
not just the unbiased subclass or any chosen classifier family.
\end{theorem}

The upper bound in (1) reuses the current separately reviewed paper proof in
\texttt{research/\allowbreak{}2026-10-27-minimax-audit/\allowbreak{}theorem/\allowbreak{}critical.tex}, which supplies
$M\le(165/56)\rho$ and $W\le1350a^{-2}L(1+\log L)$.
That source is preserved, and its theorem is not claimed as new here.
The new lower proof below is self-contained. Its finite-grid sign-testing
argument follows the established entropy mechanism of Chen--Li--Qiao
\cite[Lemma 4.1 and Appendix D]{clq}; we prove the needed Bernoulli,
separated-gap, one-sided-cost version explicitly rather than silently
substituting their Gaussian formulation.

For the population minimax problem, the candidate modulus is
$\Psi(B^{-1/2})\asymp\min\{\rho,\sqrt{\rho L}/(a\sqrt B)\}$.
Thus (1) shows that a proof through an unbiased single-coin variance--work
bound of order $\rho L$ is impossible for this fixed law. Changing only the
allocation between its variance and work cannot remove the iterated logarithm.
A matching population algorithm would need a different argument or estimator
class. No reduction from arbitrary biased or cross-coin algorithms to this
class is asserted.

\section{A finite-grid sign-testing entropy bound}
\begin{lemma}\label{lem:sign}
There are universal $c_s>0$ and an integer $K_s$ such that the following holds.
Let $0<d_K<\cdots<d_1\le1/128$, with $d_j/d_{j+1}\ge2$.
One randomized observable procedure receives iid Bernoulli flips of unknown
parameter $1/2\pm d_j$, stops, and returns a sign. Suppose its error is at most
$1/10$ at all $2K$ parameter values and its expected stopping time is finite
there. Put
\[
 c_j=d_j^2\E_{-d_j}\tau,\qquad \bar c=K^{-1}\sum_{j=1}^Kc_j.
\]
For $K\ge K_s$,
\[
 \sum_{j=1}^K d_j^2\E_{-d_j}\tau\ge c_sK\log K.          \tag{2}
\]
Correctness and termination at the fair coin are not assumed.
\end{lemma}
\begin{proof}
Set $b_0=1/1024$, $n_j=\lfloor b_0/d_j^2\rfloor$, and
$u_j=\lceil4\E_{-d_j}\tau\rceil$. Since $d_j\le1/128$,
$n_j\ge b_0/(2d_j^2)\ge1$.
The event of returning minus by time $n_j$ is measurable in the independent
seed and the first $n_j$ flips, including unused flips if the procedure stops
earlier. Between the laws $1/2-d_j$ and $1/2+d_j$, the KL divergence of these
full finite prefixes is at most
\[
 n_j\frac{(2d_j)^2}{(1/2+d_j)(1/2-d_j)}
 \le(64/3)n_jd_j^2\le1/48.
\]
Here Bernoulli KL is at most its chi-square divergence. Pinsker bounds total
variation by less than $1/8$. Since the event is an error under the plus law,
its probability under the minus law is at most $1/10+1/8$.
Markov gives $\Pp_{-d_j}(\tau>u_j)\le1/4$.
Consequently the finite-window event
\[
 A_j=\{\text{output minus and } n_j<\tau\le u_j\}
\]
has probability at least $1-1/10-(1/10+1/8)-1/4>1/4$ under the minus law.

Let $P_0$ denote the fair-coin experiment with the same independent seed and
policy. At a finite prefix of length $u_j$, its KL from the minus law is at
most $4d_j^2u_j\le16c_j+1$. For any event of probability $p\ge1/4$ under
that minus law and probability $z$ under $P_0$, binary data processing and
$\operatorname{kl}(p,z)\ge p\log(1/z)-\log2$ give
\[
 P_0(A_j)\ge\exp(-64c_j-8).                              \tag{3}
\]
All information comparisons here concern deterministic finite prefixes.
There is no optional-stopping identity under $P_0$, which may fail to halt.

At least $K/2$ indices form a set $G$ with $c_j\le2\bar c$. For those indices,
\[
 (n_j,u_j]\subset
 [b_0/(2d_j^2),(8\bar c+1)/d_j^2].
\]
The numbers $d_j^{-2}$ are separated by factors at least four. At every
finite time, the number of these intervals containing that time is therefore
at most
\[
 H(\bar c)=1+\log_4((16\bar c+2)/b_0).
\]
On every realization under the fair law, all events $A_j$ concern the same
terminal time and output. Summing their indicators, then taking expectations,
gives $\sum_{j\in G}P_0(A_j)\le H(\bar c)$, even on paths that never halt.
Together with (3),
\[
 (K/2)e^{-128\bar c-8}
 \le1+\log_4(1024(16\bar c+2)).                          \tag{4}
\]
This entails $\bar c\ge(\log K)/512$ for all sufficiently large universal
$K$. One can take $K_s=\lceil e^{32}\rceil$: if
$\bar c<(\log K)/512$ and $x=\log K\ge32$, the left side of (4) exceeds
$\tfrac12e^{3x/4-8}\ge\tfrac12e^{x/2}\ge x^2/16$,
while its right side is at most $1+\log_4(32x+2048)<6+x/4$.
These inequalities contradict $x\ge32$. Thus $c_s=1/512$ suffices.
\end{proof}

\section{What unbiasedness for every global-rarity allocation forces}
For a fixed procedure first define the nonnegative extended conditional moments
\[
 M_s(d)=\E_{d,s}Y^2,
 \quad w_s(d)=\E_{d,s}T.
\]
These are measurable nonnegative conditional moment functions, defined by the
explicit iid Bernoulli kernel and independent seed. Their integrability is
proved before signed manipulations: the admissible choices $q=0$ and
$q\equiv\rho$ imply $\int M_-d\mu\le M$ and
$\int M_+d\mu\le M/\rho$, as well as finite integrated work on both signs.
All conditional moments used below are thus finite outside one common
$\mu$-null set. On the set $\{M_s(d)<\infty\}$ define the integrable signed
conditional mean $m_s(d)=\E_{d,s}Y$, and assign it zero on its exceptional null
set. This avoids interpreting any undefined difference of infinite integrals.

Unbiasedness at $q=0$ gives $\int m_-d\mu=0$.
For every measurable $E$ of positive mass, take
$q=\min(1,\rho/\mu(E))\ind_E$. Subtracting its unbiasedness identity from
the one at zero proves $\int_E(m_+-m_-)d\mu=\mu(E)$.
Hence
\[
 m_+(d)-m_-(d)=1\quad\mu\text{-almost everywhere}.         \tag{5}
\]
This uses the global constraint only; it does not assume $q\le\rho$ pointwise.
Conditional Jensen at the constant prior $q\equiv\rho$ now gives
\[
 \E_\rho Y^2\ge\int[(1-\rho)m_-^2+\rho(m_-+1)^2]d\mu
               =\int m_-^2d\mu+\rho\ge\rho.
\]
Consequently
\[
 M\ge\rho,\qquad V\ge\rho(1-\rho),\qquad
 M\le V+\rho^2\le\tfrac43V.                             \tag{6}
\]
The response identity is stronger than a single unbiasedness statement at a
known mixture: its pointwise content will make a sign-testing reduction possible.

\section{Removing artificial thinning and controlling all product allocations}
The event $A=\{T\ge1\}$ is determined by the independent seed before any
observation. Its probability $p$ is independent of every latent state. It is
positive, since otherwise the expectation of $Y$ could not vary with $q$.
Put $c=\E[Y\ind_{A^c}]$, also independent of $q$.
Use the conditional seed law on $A$ to define the active procedure
\[
 Z=pY+c\quad\hbox{conditional on }A,
 \qquad T_Z=T\quad\hbox{conditional on }A.
\]
This is a valid mathematical single-coin procedure. Conditioning can be
implemented by resampling independent seeds until $A$ occurs, before making
any charged observation. The constants $p,c$ are properties of the original
policy, independent of $D,S,q$; the reduction may give them to the test.
The active procedure always makes at least one flip, and
\[
 \E_q Z=\theta(q),\qquad W_Z=W/p\ge1.
\]
For $p<1$, Cauchy--Schwarz on $A^c$ and completing the square give
\begin{align*}
 \E_q Z^2
 &=p\E_q[Y^2\ind_A]+2c\theta-c^2\\
 &\le pM+2c\theta-c^2/(1-p)
 \le pM+(1-p)\theta^2
 \le m:=pM+\rho^2.
\end{align*}
For $p=1$ take $c=0$ and the same final bound holds directly.
Applying (6) to $Z$ shows
$pM\ge\rho-(1-p)\rho^2\ge\rho(1-\rho)$.
Write $P=MW$. The resulting useful bounds are
\[
 \rho\le m\le\tfrac43pM\le\tfrac43P,
 \qquad mW_Z\le\tfrac43P,
 \qquad P\ge\tfrac34\rho.                               \tag{7}
\]
Thus random refusal to observe the coin cannot evade the product argument.

\section{Good gaps and the same-coin sign reduction}
Use conditional means and moments for $Z$ in this section, denoted again by
$m_s(d),M_s(d),w_s(d)$. We have $\sup_q\E_qZ^2\le m$ and response identity (5).
The set $B_-:=\{M_->1/64\}$ has mass at most $64m$.
The set
\[
 B_+:=\{M_+-M_->2m/\rho\}
\]
has mass less than $\rho$: otherwise nonatomicity of the exact power law gives
a subset $E\subset B_+$ of mass $\rho$, and the admissible $q=\ind_E$ would
give $\E_q Z^2=\int M_-d\mu+\int_E(M_+-M_-)d\mu>2m$, a contradiction.
After excluding also the common null sets above, the good set $G$ has complement
of mass at most $64m+\rho\le65m$. On $G$,
\[
 |m_-|\le1/8,\qquad m_+=m_-+1\ge7/8,
 \qquad M_-\le1/64,\qquad M_+\le1/64+2m/\rho.             \tag{8}
\]

Let $n=\lceil256(1+m/\rho)\rceil\le513m/\rho$; the last inequality uses
$m\ge\rho$. Run $n$ independent restarts of the active procedure on the
\emph{same} fixed-gap coin, with independent seeds and consecutive unused flip
blocks, and output plus iff their average exceeds $1/2$.
The test is independent of which good gap is supplied. Conditional on its
fixed latent state, each restart has the stated law, and the outputs are iid.
To justify this with unbounded stopping, condition on each finite first stopping
time and finite observed prefix; unused iid coordinates retain their original
law. Sum over the countably many prefixes, whose probabilities sum to one
because conditional expected work is finite. Repeating this argument proves
the restart assertion and additivity of expected work.
By (8), the threshold is separated from either mean by at least $3/8$.
Chebyshev gives error at most $1/10$ on both signs, because
\[
 \frac{64}{9n}\max\{1/64,1/64+2m/\rho\}<1/10.
\]
Its expected work under the minus sign at gap $d$ is exactly $nw_-(d)$.
This construction has not drawn additional latent gaps: it repeatedly observes
one fixed coin, as required by Lemma~\ref{lem:sign}.

\section{Selecting geometric gaps from the continuous law}
For integers $j\ge3$ set
\[
 I_j=(a4^{-j}/2,a4^{-j}],\qquad \mu(I_j)=\tfrac34\,16^{-j}.
\]
Their intervals are disjoint and any choices $d_j\in I_j$ satisfy
$d_j/d_{j+1}>2$; also $d_j\le a/64=1/256<1/128$.
Let
\[
 J=\left\lfloor\log_{16}\frac{3}{520m}\right\rfloor,
 \qquad K=J-2,
\]
and consider the case $J\ge3$. For $3\le j\le J$, each annulus has mass
at least $130m$, so $\mu(G\cap I_j)\ge\mu(I_j)/2$.
Choose $d_j\in G\cap I_j$ with
\[
 w_-(d_j)\le\frac4{\mu(I_j)}\int_{I_j}w_-\,d\mu.       \tag{9}
\]
Such a choice follows from the average on $G\cap I_j$. The integral is strictly
positive because the active procedure always makes at least one observation,
so $w_-(d)\ge1$. All exceptional null sets have already been removed.
These are existential selections used only in the proof; the procedure is not
given their values or the latent gap.

Since $d_j^2\le a^2 16^{-j}$, summing (9) yields
\[
 \sum_{j=3}^Jd_j^2 w_-(d_j)
 \le\frac{16a^2}{3}\int w_-\,d\mu
 \le\frac{16a^2}{3}W_Z.
\]
Apply Lemma~\ref{lem:sign} to the same test from the preceding section.
For $K\ge K_s$, its cost lower bound and (7) give
\[
 c_sK\log K
 \le\sum_{j=3}^Jd_j^2 nw_-(d_j)
 \le\frac{16a^2}{3}\frac{513m}{\rho}W_Z
 \le\frac{3648a^2P}{\rho}.                              \tag{10}
\]
The coefficient $3648=(16/3)\cdot513\cdot(4/3)$ is deliberately unoptimized.
This proves a quantitative implicit product obstruction with every parameter
and exceptional set accounted for.

\section{Completion of the product theorem and what it leaves open}
We deduce an explicit order from (10) as $\rho\downarrow0$, keeping $a=1/4$.
If $P\ge\sqrt\rho$, then $P\ge c\rho L(1+\log L)$ for all sufficiently small
$\rho$, since $\sqrt\rho L(1+\log L)\to0$.
If $P<\sqrt\rho$, (7) gives $m<(4/3)\sqrt\rho$, and the definition of $K$
implies
\[
 K\ge\frac{\tfrac12\log(1/\rho)-\log(2080/9)}{\log16}-3.
\]
For sufficiently small universal $\rho$, this is at least a universal positive
multiple of $L$, exceeds $K_s$, and has $\log K\ge\tfrac12\log L$.
Equation (10) then gives $P\ge c\rho L(1+\log L)$, with a universal $c>0$.
All thresholds here are universal because $a$ was fixed.
For the remaining compact interval of rarities, reduce $c$ if necessary and use
$P\ge3\rho/4$ from (7). Thus the raw-second-moment product lower bound holds
for every $0<\rho\le1/4$.
Finally $M\le(4/3)V$ by (6), so the same order lower bound holds for $VW$.
The previously available upper construction supplies the other direction,
proving Theorem~\ref{thm:main}.

The result explains a real statistical loss in the unbiased single-coin route:
its pointwise response identity permits recovery of a uniformly accurate sign
test at many separated gap scales. The entropy cost of those scales forces the
iterated logarithm. This is stronger than failure of the particular likelihood
dictionary or a numerical estimator search, because the proof permits arbitrary
single-coin observable policies and arbitrary separate variance--work allocations.

The result alone does not exclude biased estimators. The next section proves
the separate extension stated in Theorem~\ref{thm:bias}. It remains possible
that nonlinear use of population data or interaction between multiple genuine
coins attains the candidate
$\Psi$ rate. The current proof neither constructs such an algorithm nor proves
that it cannot exist. An extension would require a valid reduction from the
full population minimax class to one of these single-coin classes, or a new direct
information argument. Such a reduction is an open premise, not an accepted
step of this note.

\section{Proof of the small-bias extension}
Put $\eta=1/64$. Consider a procedure with
$\sup_q|\E_qY-\theta(q)|\le\eta\rho$ and the same finite separate bounds.
As before, the nonnegative conditional moments are integrated first under
$q=0$ and $q\equiv\rho$, and the signed means are defined only where square
integrable. Put $b(d)=\E_{d,-}Y$ and
$f(d)=\E_{d,+}Y-\E_{d,-}Y$. The two population bias bounds give
\[
 \int f\,d\mu=\frac{\E_\rho Y-\E_0Y}{\rho}\ge31/32.
\]
The law of total variance, conditioning first on $D$ and then on $S$, gives
\[
 V\ge\Var_\rho(Y)
 \ge\rho(1-\rho)\int f^2\,d\mu
 \ge\rho(1-\rho)(31/32)^2\ge\tfrac23\rho.                \tag{14}
\]
There is no pointwise response assumption in this step. Since every absolute
population mean is at most $65\rho/64$, also
\[
 M\le V+(65\rho/64)^2\le\tfrac32 V.                     \tag{15}
\]

Normalize the first-observation event as before. Its probability $p$ is
positive: an output taking no observations has one common expectation and
cannot approximate both $\theta=0$ and $\theta=\rho$ within $\rho/64$.
For $c=\E[Y\ind_{\{T=0\}}]$, use the active output $Z=pY+c$ conditional on
$T\ge1$. Its expectations are exactly those of $Y$, so the bias guarantee is
unchanged, and $W_Z=W/p\ge1$. The same completed-square computation, with
the actual mean $\E_qY$ in place of $\theta$, yields
\[
 \sup_q\E_qZ^2\le pM+(65\rho/64)^2\le m:=pM+2\rho^2.
\]
Applying (14) to $Z$ shows
$m\ge2\rho/3$ and $pM\ge2\rho/3-2\rho^2\ge\rho/6$.
Consequently, for $P=MW$,
\[
 m\le4pM\le4P,\qquad mW_Z\le4P,\qquad P\ge\rho/6.       \tag{16}
\]
Indeed $2\rho^2\le\rho/2\le3pM$. Thus arbitrary pre-observation thinning
and moment--work reallocations are again covered.

Use conditional moments for $Z$ in what follows. The negative-moment bad set
$\{M_->1/64\}$ has mass at most $64m$. As in the unbiased proof,
$\{M_+-M_->2m/\rho\}$ has mass less than $\rho$.
The additional response bad set
\[
 B_f=\{d:f(d)<3/4\}
\]
also has mass less than $\rho$. If it had mass at least $\rho$, choose a
measurable $E\subset B_f$ of mass $\rho$ and the admissible $q=\ind_E$.
Then
\[
 \E_{\ind_E}Z-\E_0Z=\int_E f\,d\mu<3\rho/4,
\]
while the two bias bounds force this difference to be at least
$\rho-2\rho/64=31\rho/32$, a contradiction. This argument uses only
global-rarity alternatives and does not assume pointwise classification errors.

After excluding these three bad sets and the common null sets, the good set
has complement of mass at most $64m+2\rho\le67m$. On it,
\[
 |m_-|\le1/8,\qquad m_+\ge5/8,\qquad
 M_+\le1/64+2m/\rho,\qquad M_-\le1/64.
\]
Restart $Z$ independently
$n=\lceil512(1+m/\rho)\rceil\le1282m/\rho$ times on the same fixed coin;
the inequality uses $m/\rho\ge2/3$. Threshold their average at $3/8$.
Both mean separations are at least $1/4$, so Chebyshev bounds both errors by
\[
 \frac{16(1/64+2m/\rho)}{n}<1/10.
\]
The iid unused-suffix argument already proved gives expected work $nw_-(d)$
on a negative gap. No additional genuine coin or magnitude is introduced.

Use the same annuli $I_j$, now retaining
\[
 j\le J=\left\lfloor\log_{16}\frac3{536m}\right\rfloor,
 \qquad K=J-2.
\]
Each retained annulus has mass at least $134m$, and the good set occupies at
least half. The identical averaging selection and sign-entropy lemma give,
for $K\ge K_s$,
\[
 c_sK\log K
 \le\sum_jd_j^2nw_-(d_j)
 \le\frac{16a^2}{3}\frac{1282m}{\rho}W_Z
 \le\frac{30000a^2P}{\rho}.                             \tag{17}
\]
Here $30000>(16/3)\cdot1282\cdot4$ is a conservative numerical constant.
If $P<\sqrt\rho$, (16) shows $m\le4\sqrt\rho$, so $K$ is at least a
universal positive multiple of $\log(1/\rho)$ for sufficiently small $\rho$.
Equation (17) proves the required order lower bound. If $P\ge\sqrt\rho$,
that order already holds for small $\rho$; extend to the remaining compact
interval using $P\ge\rho/6$. Finally (15) transfers the lower bound from
$MW$ to $VW$. The unbiased upper construction also belongs to this bias class,
which proves both sides of Theorem~\ref{thm:bias}.

\paragraph{Exact averaging consequence.}
For an unprojected mean of a fixed number $n$ of iid outputs of one such
procedure, a uniform expected budget $nW\le B$ implies worst-case variance
at least $V/n\ge VW/B$. Independent Poisson replication with mean count
$\lambda$ and output normalized by $\lambda$ has variance
$\E_qY^2/\lambda$, giving the corresponding raw-moment product obstruction
when $\lambda W\le B$. In particular small uniform per-coin bias does not
remove this loss from either of these averaging arguments.
These consequences do not include arbitrary nonlinear aggregation,
data-dependent repetition counts, clipping, or interaction between genuine
coins. No all-algorithm minimax conclusion follows without an additional
valid reduction.

% Appendix numbering supplied by the wrapper.
\section{Self-contained upper checkpoint and provenance}
This appendix records the existing upper argument needed for the second
inequality in (1), so the characterization does not depend on an inaccessible
working-tree file. The inspected source was
\texttt{research/\allowbreak{}2026-10-27-minimax-audit/\allowbreak{}theorem/\allowbreak{}critical.tex}, SHA-256
\begin{center}\small\texttt{0e8b9eaf762b3b09fa0d92ce1a036c2c}\\\texttt{eea8c12d00a5eea8286b852f36a81536}\end{center}
It was read without modifying it. The following is a compact restatement of
that pre-existing proof, not a new construction or discovery claim.
The source is now committed at
\par\noindent\texttt{a8a8dc49115aabdd1fba8d53d3390eb0a5eee864}.

Put $H=1+\log L$. For $j\ge0$ choose
\[
 \varepsilon_j=16^{-(j+1)},\qquad
 N_j=\lceil(a^2\rho\varepsilon_j)^{-1}\rceil,\qquad
 K_j=\lceil\log_2(a\sqrt{N_j})\rceil.
\]
For $0\le k\le K_j$, let $r_k=a2^{-(k+1)}$ and, after $n$ observations
with $h_n$ heads, define
\[
 L^+_{k,n}=(1+2r_k)^{h_n}(1-2r_k)^{n-h_n},\qquad
 L^-_{k,n}=(1-2r_k)^{h_n}(1+2r_k)^{n-h_n}.
\]
Before $N_j$, stop with output zero if a negative likelihood reaches
$(K_j+1)/\varepsilon_j$, or with output one if a positive likelihood reaches
$(K_j+1)/(\rho\varepsilon_j)$. Fix zero priority for ties. At $N_j$, if no
earlier decision was made, return the strict majority bit. Denote this bounded,
observable classifier by $Z_j$ and its work by $T_j$.

For the wrong-sign likelihood, each next multiplier has mean $1-4r_kd\le1$.
Stopping this nonnegative supermartingale at a crossing and a deterministic
horizon, and then increasing the horizon, bounds the probability of crossing
level $b$ by $1/b$. Summing over $K_j+1$ references gives wrong-crossing
probabilities at most $\varepsilon_j$ on a positive coin and
$\rho\varepsilon_j$ on a negative coin. The wrong forced-majority event is
a subset of the full-tape event, of probability at most $e^{-2N_jd^2}$.
This is Hoeffding's bound: the centered Bernoulli log moment-generating
function has second derivative at most $1/4$, so is at most $\lambda^2/8$;
use independence and exponential Markov with $\lambda=\pm4d$.
For the exact power law,
\[
 \E_\mu e^{-2ND^2}=\frac{1-e^{-2Na^2}}{2Na^2}\le\frac1{2Na^2}.
\]
No independence of a stopping event and this full-tape majority event is used.
It follows, uniformly over the global-rarity class, that
\[
 \Pp(Z_j\ne S)\le\tfrac52\rho\varepsilon_j,
 \qquad \Pp(Z_j=1,S=0)\le\tfrac32\rho\varepsilon_j.       \tag{11}
\]

For $d\ge N_j^{-1/2}$, choose $k=\lceil\log_2(a/d)\rceil\le K_j$.
Then $d/4\le r_k\le d/2$, and the correct-sign log likelihood has mean
increment at least $2d^2/3$: use
\[
 \tfrac12\log(1-4r_k^2)
       +d\log\frac{1+2r_k}{1-2r_k}
 \ge-\tfrac83r_k^2+4dr_k\ge\tfrac23d^2.
\]
Its upward increment is less than one. The bounded stopped-sum identity for
the crossing time $\sigma$ at a finite horizon $n$ gives
$g\E(\sigma\wedge n)\le A_{j,s}$, where
\[
 A_{j,s}=1+\log(K_j+1)+\log(1/\varepsilon_j)+s\log(1/\rho).
\]
Monotone convergence proves finite expected crossing time and, including the
cap and the trivial small-gap case,
\[
 \E(T_j\mid d,s)\le\min\{N_j,(3/2)A_{j,s}/d^2\}.        \tag{12}
\]
Write $c_0=3/2$. Direct integration of the negative contribution gives
\[
 \E\min\{N,c_0A/D^2\}
 \le(c_0A/a^2)[1+\log_+(Na^2)].
\]
For positives, layer cake and
$\E[q\ind_E]\le\min\{\rho,\mu(E)\}$ give
\[
 \E[q\min\{N,c_0A/D^2\}]
 \le\int_0^\rho\min\{N,c_0A/(a^2u)\}\,du
 \le(c_0A/a^2)[1+\log_+(Na^2\rho)].
\]
The exact integral is $N\rho$ when $Na^2\rho\le c_0A$, and otherwise is
$(c_0A/a^2)[1+\log(Na^2\rho/(c_0A))]$. Thus this step also enforces the
global, rather than pointwise, rarity condition.

Elementary bounds from the definitions are
\[
\begin{gathered}
 N_ja^2\le2/(\rho\varepsilon_j),\qquad K_j+1\le5(L+j+1),\\
 A_{j,0}\le7H(j+1),\qquad A_{j,1}\le7H(j+1)+L,\\
 1+\log(N_ja^2)\le5L(j+1),\qquad
 1+\log(N_ja^2\rho)\le5(j+1).
\end{gathered}
\]
For example, $\log(L+j+1)\le\log L+j+1$ and $4\log2<3$ prove the
bound on $A_{j,0}$. Substitution in (12) and the two integrals gives
\[
 \E T_j\le(225/2)a^{-2}LH(j+1)^2.                        \tag{13}
\]

Independently choose a level $J$ with $p_j=2^{-(j+1)}$. For $J=0$ return
$Y=Z_0/p_0$ on one fresh coin. For $J=j\ge1$, run $Z_j$ and $Z_{j-1}$
on successive unused blocks of that same coin, set $\Delta_j=Z_j-Z_{j-1}$,
and return $Y=\Delta_j/p_j$. The bounded stopping times make the iid suffix
argument elementary, as in the restart argument above; both decisions share
the same latent sign. A disagreement implies an error of one classifier.
Hence (11) gives
\[
 \E Z_0^2\le\tfrac{35}{32}\rho,
 \qquad \E\Delta_j^2\le\tfrac{85}{32}\rho16^{-j}\quad(j\ge1).
\]
Tonelli and the exact geometric sum now give
\[
 \E Y^2\le\frac{35}{16}\rho+
                 \frac{85}{16}\rho\sum_{j\ge1}8^{-j}
            =\frac{165}{56}\rho.
\]
The disagreement bounds also give absolute summability of expected increments,
so signed expectations may be telescoped after integrability is established:
$\E Y=\lim_j\E Z_j=\theta$. Work costs at most twice (13), and
$\sum_jp_j(j+1)^2=6$, yielding
\[
 \E T_Y\le1350a^{-2}LH.
\]
All operations are on one observable coin with independent seed; the random
level is unbounded but the separate expectations are finite. This completes
the upper checkpoint used in (1).

\begin{thebibliography}{9}
\bibitem{clq} L. Chen, J. Li, and M. Qiao.
\emph{Towards Instance Optimal Bounds for Best Arm Identification}.
COLT 2017; arXiv:1608.06031v2.
\url{https://arxiv.org/abs/1608.06031}.
\bibitem{critical} Current workspace continuation,
\emph{A critical-margin upper bound with a remaining iterated logarithm},
\texttt{research/\allowbreak{}2026-10-27-minimax-audit/\allowbreak{}theorem/\allowbreak{}critical.tex}.
Reviewed paper proof, not Lean-checked; preserved without modification.
\end{thebibliography}
